\section{Positive approximation below the bound \texorpdfstring{$2d-1$}{2d-1}}
\label{human:coherent-section}

This section proves the first branch directly. The approximating measures
are positive probabilities throughout, and converge in first norm to the
actual pinned distance law. The only specialized external input is the
following published radial-projection theorem.

\begin{theorem}[Orponen's averaged radial-projection theorem]
\label{human:radial}
Let \(\alpha>1\). For two compactly supported planar probabilities
\(\mu,\nu\) with disjoint supports and bounds
\[
 \mu(B(z,r))+\nu(B(z,r))\le C r^\alpha,
 \qquad 0<r\le1,
\]
there is a number \(q=q(\alpha)>1\) such that, for
\(\mu\)-almost every \(x\),
\begin{equation}\label{human:radial-moment}
 (\Theta_x)_*\nu=\rho_x\sigma,
 \qquad \Theta_x(y)=\frac{x-y}{|x-y|},
 \qquad
 B:=\int\|\rho_x\|_{L^q(\sigma)}^q\,d\mu(x)<\infty.
\end{equation}
Here \(\sigma\) is normalized arclength on the circle.
\end{theorem}

This is the form of Orponen's result \cite{orponen} stated in
\cite[Theorem 3.7]{giow}. A common translation and dilation puts the
supports in the unit disk required there. Reversing the radial direction
does not change an angular norm. In particular, different Frostman
exponents \(s,t>1\) are allowed: choose
\(1<\alpha<\min\{s,t\}\).

\subsection{A projection estimate that permits correlated angles}

We give the analytic estimate in detail, since the two directions used
later are functions of the same pin. The Fourier convention is
\(\widehat\lambda(\xi)=\int e^{-2\pi i x\cdot\xi}\,d\lambda(x)\).

\begin{lemma}\label{human:bandlemma}
Fix \(1<s<2\) and geometric constants. Let \(0<r\le1\), and let
\(\lambda\) be a probability supported in \(B(0,C_0r)\), with
\(\lambda(B(z,t))\le Kt^s\) for all \(z,t>0\).
Let \(F(\theta)\) be the density of the orthogonal projection
\(x\mapsto x\cdot\theta\), defined for almost every \(\theta\).
Suppose a finite positive measure \(\kappa\) carries measurable
directions \(\theta,\phi\) and a scalar shift \(v\), with both
angular marginals bounded by \(A\sigma\), and
\[
 d(\theta,\phi)\le C_1r,\qquad |v|\le C_1r^2.
\]
Then
\begin{equation}\label{human:bandbound}
 \int\|F(\theta)(\cdot-v)-F(\phi)\|_2^2\,d\kappa
 \le C AKr^{2s-2}.
\end{equation}
No independence of directions or shifts is assumed.
\end{lemma}

\begin{proof}
For a frequency radius \(R>0\), the Gaussian Fourier identity gives
\begin{align}
 E_R&:=\int e^{-|\xi|^2/R^2}|\widehat\lambda(\xi)|^2\,d\xi
 \label{human:gaussian}\\
 &=cR^2\iint e^{-c'R^2|x-x'|^2}\,d\lambda(x)d\lambda(x')
 \le C_sKR^{2-s}.\nonumber
\end{align}
The last bound follows by the Frostman estimate on the ball of radius
\(R^{-1}\) and on the successive doubled annuli. Positivity of the
spatial Gaussian also proves, for the signed coordinate measures,
\begin{equation}\label{human:coordinate}
 \int e^{-|\xi|^2/R^2}|\widehat{x_k\lambda}(\xi)|^2\,d\xi
 \le C r^2E_R,\qquad k=1,2:
\end{equation}
in the spatial double integral, the additional factor has absolute
value \(|x_kx'_k|\le C_0^2r^2\).

Choose smooth scalar-frequency cutoffs \(\psi_R\), \(R\in2^{\mathbb Z}\),
supported where \(R/2\le|\tau|\le2R\), so that
\(\sum_R\psi_R(\tau)^2=1\) for \(\tau\ne0\). Put
\[
 U_R(\theta)=\mathcal F^{-1}_\tau
       [\psi_R(\tau)\widehat\lambda(\tau\theta)].
\]
Polar coordinates and \eqref{human:gaussian} give
\begin{equation}\label{human:band-energies}
 \int\|U_R(\theta)\|_2^2\,d\sigma(\theta)\le CKR^{1-s},
 \qquad
 \int\|\partial_\theta U_R(\theta)\|_2^2\,d\sigma(\theta)
 \le CKr^2R^{3-s}.
\end{equation}
For the derivative bound use \eqref{human:coordinate}: differentiating
\(\widehat\lambda(\tau\theta)\) inserts a bounded linear combination
of \(\tau\widehat{x_k\lambda}(\tau\theta)\).
The projected densities exist and are jointly square integrable in
angle and scalar variable. Indeed, summing the first estimate over
\(R\ge1\) converges because \(s>1\); for \(R<1\) use
\(|\widehat\lambda|\le1\), giving instead the summable bound \(CR\).
Plancherel then identifies the projected measures with these densities
outside one angular null set.

The fundamental theorem of calculus for the smooth Hilbert-valued
function \(U_R\), along the shorter angular arc, gives
\[
 \|U_R(\theta)-U_R(\phi)\|_2^2
 \le Cr\int_{d(z,\theta)\le C_1r}
                     \|\partial_zU_R(z)\|_2^2\,dz.
\]
Integration with the first marginal and Fubini contribute another
factor \(Ar\). The trivial squared-difference bound uses both
marginals. Together with \eqref{human:band-energies} these yield
\[
 \int\|U_R(\theta)-U_R(\phi)\|_2^2\,d\kappa
 \le CAK\min\{1,(r^2R)^2\}R^{1-s}.
\]
For translations, the pointwise scalar Fourier multiplier bound
\[
 |e^{-2\pi i\tau v}-1|^2
 \le C\min\{1,(r^2R)^2\},\qquad |\tau|\asymp R,
\]
gives the same estimate. It is valid before integrating in the pin
parameter, so the shift may be correlated with either direction.

For each parameter outside the angular marginal null sets, Plancherel
and the squared partition of unity express the squared norm of the
full difference as the sum of its band squared norms. The squared
triangle inequality and Tonelli therefore reduce the assertion to
\[
 \sum_{R\in2^{\mathbb Z}}\min\{1,(r^2R)^2\}R^{1-s}
 \le C_s r^{2s-2}.
\]
Below \(R=r^{-2}\) the summand is \(r^4R^{3-s}\); above it the
summand is \(R^{1-s}\). Both geometric sums have the stated bound.
\end{proof}

\subsection{Convergence of one fixed family of positive approximations}

\begin{proposition}\label{human:coherent-measure}
Let \(\mu,\nu\) be compactly supported planar probabilities with
positive separation. Suppose \(\mu\) is \(s\)-Frostman, where
\(1<s<2\), and \(\nu\) has a Frostman exponent greater than one.
If the number \(N_n\) of occupied source dyadic squares of side
\(2^{-n}\) satisfies
\begin{equation}\label{human:covering-bound}
 N_n\le C2^{un},\qquad u<2s-1,
\end{equation}
then \((d_y)_*\mu\) is absolutely continuous for
\(\nu\)-almost every \(y\).
\end{proposition}

\begin{proof}
Apply Theorem~\ref{human:radial}. Write
\(M(x)=\|\rho_x\|_q^q\), so \(\int M\,d\mu=B<\infty\).
For each positive-mass half-open dyadic square \(Q\), put
\(p_Q=\mu(Q)\) and \(\mu_Q=p_Q^{-1}\mu|_Q\). Choose once and
for all a point \(x_Q\in Q\cap\operatorname{supp}\mu_Q\) where the
radial density exists and
\begin{equation}\label{human:center-choice}
 M(x_Q)\le\frac2{p_Q}\int_Q M\,d\mu.
\end{equation}
Such points exist by averaging. There are only countably many nodes,
and at every level
\begin{equation}\label{human:center-sum}
 \sum_{Q\text{ at level }n}p_QM(x_Q)\le2B.
\end{equation}

For \(x\in Q\), define
\[
 L_Q(x,y)=\frac{x_Q-y}{|x_Q-y|}\cdot(x-y).
\]
Let \(f_n\) be the density of the pushforward of \(\mu\times\nu\)
by \((x,y)\mapsto(y,L_{Q_n(x)}(x,y))\), relative to
\(\nu(dy)\,dt\). These densities exist: each conditional source has
finite first energy, its orthogonal projections have square-integrable
densities for almost every direction, and the chosen center's radial
pin law is absolutely continuous. The countably many conditional
densities can be chosen jointly measurably. Thus
\[
 f_n\ge0,\qquad\int f_n\,d\nu\,dt=1.
\]
Positive separation and Taylor's theorem give, uniformly for source
and pin points and sufficiently large \(n\),
\begin{equation}\label{human:affine-error}
 |L_{Q_n(x)}(x,y)-|x-y||\le C2^{-2n}.
\end{equation}

We compare consecutive levels. Put \(r=2^{-n}\), let \(Q\) be a
child of \(P\), and apply both affine maps to the same \(\mu_Q\).
Recenter this source at \(b=x_Q\). With
\(\theta=\Theta_{x_P}(y)\) and \(\phi=\Theta_b(y)\),
\[
 L_P(x,y)=\theta\cdot(x-b)+\theta\cdot(b-y),\qquad
 L_Q(x,y)=\phi\cdot(x-b)+|b-y|.
\]
The directions differ by \(O(r)\), while the constants differ by
\begin{equation}\label{human:quadratic-shift}
 |\theta\cdot(b-y)-|b-y||
 =|b-y|\,|\theta\cdot\phi-1|\le Cr^2.
\end{equation}
The recentering is what gives this quadratic error.

Retain only pins where both selected angular densities are at most
\(A\ge1\). The two retained angular marginals are bounded by
\(A\sigma\), and the discarded pin mass is at most
\[
 A^{1-q}\bigl(M(x_P)+M(x_Q)\bigr).
\]
The translated conditional source is supported in a ball of radius
\(Cr\) and has Frostman constant at most \(C_\mu/p_Q\).
Lemma~\ref{human:bandlemma} gives retained integrated squared
second-norm error at most
\[
 \frac{CA}{p_Q}r^{2s-2}.
\]
For each pin, the two conditional densities occupy the union of two
intervals of total length \(Cr\). Cauchy--Schwarz, first on that
union and then over retained pins, bounds their integrated first-norm
difference by \(CA^{1/2}p_Q^{-1/2}r^{s-1/2}\).
On discarded pins use the bound two for the difference of probability
densities.

Multiply by \(p_Q\) and sum. The retained terms satisfy
\(\sum_Q\sqrt{p_Q}\le\sqrt{N_{n+1}}\); the discarded parent
terms sum with the parent masses, because the child masses sum to
their parent's mass. Equation~\eqref{human:center-sum} therefore gives
\begin{equation}\label{human:coherent-comparison}
 \|f_{n+1}-f_n\|_{L^1(\nu\times dt)}
 \le C A^{1/2}r^{s-1/2}\sqrt{N_{n+1}}
       +CB A^{1-q}
 \le C\sqrt A\,r^{(2s-1-u)/2}+CB A^{1-q}.
\end{equation}
No new centers are chosen for this comparison.

Write \(\kappa=2s-1-u>0\) and choose \(A=r^{-\kappa/2}\).
The last bound is
\[
 C r^{\kappa/4}+CB r^{\kappa(q-1)/2},
\]
which is summable over dyadic scales. Hence the same positive
probabilities \(f_n\) converge in first norm to a nonnegative density
of mass one. Equation~\eqref{human:affine-error} identifies their weak
limit as the pushforward of \(\mu\times\nu\) by
\((x,y)\mapsto(y,|x-y|)\). The first-norm limit is this same measure.
Uniqueness of disintegration gives its density on almost every pin
fiber, proving the proposition.
\end{proof}

This proves Proposition~\ref{human:coherent-interface}: a bounded
\(t\)-Frostman probability is also \(s\)-Frostman when \(s<t\), and
an \(O(r^{-u})\) ball-covering bound gives \(O(2^{un})\) occupied dyadic
squares by bounded overlap at comparable scales. The compact reduction
needed for the Borel theorem has already been proved in
Lemma~\ref{human:compact}.
